\chapter {Desenvolvimento}

Neste capítulo são detalhadas as opções técnicas tomadas durante a implementação da solução, focando-se na arquitetura do sistema, no processamento dos ficheiros bancários e na construção da interface interativa.

\section {Arquitetura da Solução}
A solução foi desenvolvida de forma modular, separando a lógica de leitura de dados da lógica de visualização. O fluxo de dados ocorre da seguinte forma:
\begin{enumerate}
    \item Leitura dos ficheiros \texttt{.ps2} através do módulo \texttt{parser.py};
    \item Conversão e limpeza dos dados utilizando a biblioteca \texttt{pandas};
    \item Renderização da interface web e gestão de eventos através do \texttt{Shiny}.
\end{enumerate}

\subsection{Módulo de Leitura (Parser)}
O módulo \texttt{parser.py} é responsável pela extração de informação bruta. Um dos desafios superados foi a interpretação do identificador único do cliente.  
Conforme implementado, o sistema divide o campo de 30 caracteres nos seus componentes lógicos:

\begin{verbatim}
campo_identificacao = linha[11:41] # 30 caracteres totais
c_nif = campo_identificacao[:9]    # Primeiros 9 sao o NIF
c_nib = campo_identificacao[9:]    # Restantes 21 sao o NIB
\end{verbatim}

\section{Engenharia de Dados}
Para garantir a consistência temporal, foi implementado um mecanismo de validação que prioriza a data presente no nome do ficheiro (por exemplo, \verb|DD_2025_01.ps2|) em detrimento de datas internas incorretas.  
Adicionalmente, são filtrados registos anteriores a 2024, de forma a evitar inconsistências gráficas e estatísticas.

\section{Interface e Interatividade}

A interface gráfica foi construída com a framework \textit{Shiny for Python}, aproveitando a sua capacidade de criar aplicações web robustas diretamente em Python. Esta escolha permitiu uma integração nativa com as ferramentas de análise de dados previamente desenvolvidas.

\vspace{0.5cm}

A interatividade é gerida através de funções reativas (\verb|@reactive.calc|), que filtram o \textit{DataFrame} global sempre que o utilizador altera o intervalo de datas ou seleciona um cliente específico. Este paradigma garante que o processamento só ocorre quando há uma mudança efetiva nos inputs.

\vspace{0.5cm}

Estas alterações de estado atualizam instantaneamente os componentes visuais, nomeadamente os gráficos gerados pela biblioteca \texttt{plotly}. A natureza dinâmica destes gráficos permite ao utilizador explorar pontos de dados individuais e exportar visualizações em tempo real.









